%% 13_sampling.tex -- Chapter 13: Sampling and Quantisation
%% Status: OUTLINE (build stage 2). Replace the outline with the full chapter when it is written.
\chapter{Sampling and Quantisation}\label{ch:sampling}
\chapterlead{An ADC does two things to a biosignal: it looks at it only at discrete instants, and it rounds every look to one of $2^N$ levels.}

\begin{specnote}[Chapter status: outline]
This chapter is an outline. It will be written in build stage 2 (the Python-only converter core: Chapters 13, 14 and 19 with Labs 6 and 9). The topics, the design question and the labs below are fixed; the derivations, worked examples and figures are still to come.
\end{specnote}

\begin{objectives}
\begin{itemize}
  \item apply the Nyquist theorem to ExG signals and predict what aliases where;
  \item derive the quantisation-noise model, its $6.02N + 1.76$\,dB result and its limits;
  \item use dither and oversampling to trade rate for resolution;
  \item estimate the jitter limit and see why it is minor at biosignal rates.
\end{itemize}
\end{objectives}

\begin{insight}[Design question]
What sample rate and how many bits does each \OB{} mode need, before any circuit is chosen?
\end{insight}

\section{Sampling and aliasing with ExG examples}\label{sec:ch13-1}
Mains harmonics and out-of-band EMG folding into the ECG band; anti-aliasing requirements set by real interferers rather than by the signal.

\begin{keyeq}[Quantisation-limited SNR and effective number of bits]
\begin{equation}
  \mathrm{SQNR} = 6.02N + 1.76\ \mathrm{dB}, \qquad \ENOB = \frac{\mathrm{SNDR} - 1.76}{6.02}
  \label{eq:sqnr}
\end{equation}
\end{keyeq}

\section{The quantisation-noise model}\label{sec:ch13-2}
Uniform error of $\pm\,\mathrm{LSB}/2$, its $\mathrm{LSB}^2/12$ power, and when the white-noise model fails (low-amplitude and DC signals).

\section{Dither and oversampling}\label{sec:ch13-3}
Adding noise to linearise the quantiser; each doubling of the sample rate buys \qty{3}{\deci\bel} of SQNR, setting up the noise shaping of \cref{ch:dstheory}.

\section{Clock jitter}\label{sec:ch13-4}
$\SNR_\mathrm{jitter} = -20\log_{10}(2\pi f_\mathrm{in}\sigma_t)$; at a few kilohertz, even nanoseconds of jitter are harmless, which is why biomedical ADCs worry about energy instead.

\begin{tryit}[Labs for this chapter]
Lab~6 (\cref{sec:lab6}) quantises sine waves, applies coherent sampling and windows, and builds the \code{adctest.py} module every later converter lab uses.
\end{tryit}

\paragraph{Read with this chapter} Pelgrom~\cite{pelgrom}, Schreier, Temes and Pavan~\cite{schreier}.
